% BHU PYQ -- Auto-generated & error-corrected
\documentclass[11pt,a4paper]{article}

\usepackage[a4paper,top=2.5cm,bottom=2.5cm,left=2.8cm,right=2.8cm,headheight=14pt]{geometry}
\usepackage{amsmath,amssymb}
\usepackage[T1]{fontenc}
\usepackage[utf8]{inputenc}
\usepackage{lmodern}
\usepackage{microtype}
\usepackage{fancyhdr}
\usepackage{enumitem}
\usepackage{hyperref}
\usepackage{lastpage}
\usepackage{parskip}

\hypersetup{colorlinks=true,urlcolor=black,linkcolor=black,
  pdftitle={BSC-08A: Descriptive Statistics and Probability Distributions -- BHU PYQ},pdfauthor={Banaras Hindu University}}

\pagestyle{fancy}\fancyhf{}
\renewcommand{\headrulewidth}{0.4pt}\renewcommand{\footrulewidth}{0.4pt}
\fancyhead[L]{\small\textit{Banaras Hindu University}}
\fancyhead[R]{\small\textit{BSC-08A: Descriptive Stats \& Prob Dist}}
\fancyfoot[C]{\small Page \thepage\ of \pageref{LastPage}}

\newlist{parts}{enumerate}{2}
\setlist[parts,1]{label=(\alph*),leftmargin=*,itemsep=5pt,topsep=3pt}
\setlist[parts,2]{label=(\roman*),leftmargin=*,itemsep=3pt,topsep=2pt}
\newcommand{\pts}[1]{\hfill\textit{[#1]}}

\begin{document}

\begin{center}
  {\large\bfseries BANARAS HINDU UNIVERSITY}\\[2pt]
  {\normalsize B.Sc. (Hons.) Semester III Examination, 2016-17}\\[6pt]
  \rule{0.8\linewidth}{0.5pt}\\[6pt]
  {\large\bfseries Paper No. BSC-08A: Descriptive Statistics and Probability Distributions}\\[4pt]
  {\normalsize Time: 3 Hours\hfill Maximum Marks: 70}
\end{center}
\rule{\linewidth}{0.4pt}
\smallskip
\noindent\textit{Instructions: Answer \textbf{five} questions including Question~1 which is compulsory. Symbols have their usual meanings.}
\bigskip

\medskip\noindent\textbf{Question 1.}
\begin{parts}
  \item What are the general rules of forming a frequency distribution? With the help of examples, explain the terms: (i) 'less than type' and 'greater than type' cumulative frequencies, (ii) width of a class interval, and (iii) mid-value of a class. \pts{6}
  \item Represent the following data by means of a histogram and cumulative frequency curve (ogive):
  \begin{center}
  \begin{tabular}{|l|c|c|c|c|c|c|c|}
  \hline
  \textbf{Weekly Wages (Rs)} & 10--15 & 15--20 & 20--25 & 25--30 & 30--40 & 40--60 & 60--80 \\ \hline
  \textbf{No. of Workers}    & 8      & 19     & 27     & 15     & 12     & 12     & 8      \\ \hline
  \end{tabular}
  \end{center} \pts{8}
\end{parts}
\medskip\rule{\linewidth}{0.2pt}

\medskip\noindent\textbf{Question 2.}
\begin{parts}
  \item What do you mean by central tendency of the data? Define mean, median, and mode as measures of central tendency. \pts{7}
  \item Determine the missing frequency in the following incomplete distribution when its mode is 35:
  \begin{center}
  \begin{tabular}{|l|c|c|c|c|c|c|c|}
  \hline
  \textbf{Class Interval} & 0--10 & 10--20 & 20--30 & 30--40 & 40--50 & 50--60 & 60--70 \\ \hline
  \textbf{Frequency}      & 10    & 12     & ?      & 20     & 14     & 12     & 10     \\ \hline
  \end{tabular}
  \end{center} \pts{7}
\end{parts}
\medskip\rule{\linewidth}{0.2pt}

\medskip\noindent\textbf{Question 3.}
\begin{parts}
  \item What do you mean by dispersion? Explain mean deviation and standard deviation with their relative merits and demerits. \pts{7}
  \item Define the coefficient of variation. The number of employees, average wages per employee, and variance of wages for two factories are given below. Find in which factory is there greater variation in the distribution of wages per employee:
  \begin{center}
  \begin{tabular}{|l|cc|}
  \hline
  & \textbf{Factory A} & \textbf{Factory B} \\ \hline
  \textbf{No. of Employees}     & 100                & 200                \\
  \textbf{Average Wages (Rs)}   & 120                & 80                 \\
  \textbf{Variance of Wages}    & 16                 & 25                 \\ \hline
  \end{tabular}
  \end{center} \pts{7}
\end{parts}
\medskip\rule{\linewidth}{0.2pt}

\medskip\noindent\textbf{Question 4.}
\begin{parts}
  \item Define raw and central moments for grouped data. Calculate the first four central moments from the following set of raw data values:
  \[
  2, \quad 4, \quad 6, \quad 8, \quad 10
  \] \pts{7}
  \item What do you mean by skewness and kurtosis? The first four central moments of a set of data were found to be $\mu_1 = 0$, $\mu_2 = 9$, $\mu_3 = -4$, and $\mu_4 = 122$ respectively. What would you infer about the skewness and kurtosis of this data? \pts{7}
\end{parts}
\medskip\rule{\linewidth}{0.2pt}

\medskip\noindent\textbf{Question 5.}
\begin{parts}
  \item Define Karl Pearson's correlation coefficient and mention its main properties. \pts{7}
  \item Explain Spearman's rank correlation coefficient. Calculate Spearman's coefficient of correlation from the following data:
  \begin{center}
  \begin{tabular}{|l|cccccccc|}
  \hline
  \textbf{Rank by X} & 1 & 5 & 6 & 2 & 3 & 7 & 4 & 8 \\
  \textbf{Rank by Y} & 3 & 7 & 8 & 1 & 2 & 6 & 5 & 4 \\ \hline
  \end{tabular}
  \end{center} \pts{7}
\end{parts}
\medskip\rule{\linewidth}{0.2pt}

\medskip\noindent\textbf{Question 6.}
\begin{parts}
  \item Explain the use of scatter diagrams in studying the correlation between two variables. Plot the scatter diagram for the following data and comment:
  \begin{center}
  \begin{tabular}{|l|cccccc|}
  \hline
  \textbf{X} & 1  & 2  & 3  & 4  & 5  & 6  \\
  \textbf{Y} & 10 & 20 & 30 & 40 & 50 & 60 \\ \hline
  \end{tabular}
  \end{center} \pts{7}
  \item What do you mean by regression? Write down the expression for the two regression lines and show how the regression coefficients are related to the correlation coefficient. State their main properties. \pts{7}
\end{parts}
\medskip\rule{\linewidth}{0.2pt}

\medskip\noindent\textbf{Question 7.}
\begin{parts}
  \item A six-faced ordinary die is thrown once. Find the probability of getting an outcome of: (i) less than 3 spots, and (ii) more than 2 spots. \pts{7}
  \item State and prove the additive and multiplicative theorems of probability. \pts{7}
\end{parts}
\medskip\rule{\linewidth}{0.2pt}

\medskip\noindent\textbf{Question 8.}
Write short notes on any \textbf{three} of the following:
\begin{parts}
  \item Bar and pie diagrams
  \item Normal distribution and its properties
  \item Bayes' theorem
  \item Probability mass function (p.m.f.) and probability density function (p.d.f.)
\end{parts}
\pts{14}

\vfill
\rule{\linewidth}{0.4pt}
\begin{center}\small
  \href{https://bhu-exam.vercel.app/}{Click here to visit our website}\\[4pt]
  \href{https://me-aryan.vercel.app/}{Click here to visit our contributor}\\[4pt]
  \href{https://quashan-me.vercel.app/}{Click here to visit our contributor}
\end{center}

\end{document}